% Style configuration

%%%%%%%%%%%
% Fonts
%%%%%%%%%%%
% We use Linux Libertine as main font
\KOMAoptions{DIV=last} % recalculate margins
\setmainfont{Linux Libertine O}
\setsansfont{Linux Biolinum O}
\setmonofont[Scale=0.875]{Droid Sans Mono}

% Abbreviations (from "foreign" package) are not italicized
\renewcommand{\foreignabbrfont}{}

\emergencystretch=1em % Try to avoid over-full boxes

% Table of contents
\setcounter{tocdepth}{2}
% Section numbers
\setcounter{secnumdepth}{3}


%%%%%%%%%%%
% Colours
%%%%%%%%%%%

% Custom color scheme

% Based on solarized
\definecolor{CSblack}{RGB}{0,0,0}
\definecolor{CSgrey}{HTML}{839496}
\definecolor{CSlightgrey}{HTML}{93A1A1}
\definecolor{CSyellow}{HTML}{B58900}
\definecolor{CSorange}{HTML}{CB4B16}
\definecolor{CSred}{HTML}{DC322F}
\definecolor{CSmagenta}{HTML}{D33682}
\definecolor{CSviolet}{HTML}{6C71C4}
\definecolor{CSblue}{HTML}{268BD2}
\colorlet{CSlightblue}{CSblue!75}
\definecolor{CScyan}{HTML}{2AA198}
\definecolor{CSgreen}{HTML}{859900}
 % Set colour scheme
% Colours used throughout the document

% Specify additional colours
\definecolor{darkblue}{RGB}{0,0,128}
\definecolor{lightblue}{RGB}{224,224,255}
\definecolor{darkred}{RGB}{128,0,0}
\definecolor{lightred}{RGB}{255,224,224}
%\definecolor{darkgreen}{RGB}{0,128,0}
%\definecolor{lightgreen}{RGB}{224,255,224}
\definecolor{darkpurple}{RGB}{128,0,128}
\definecolor{lightpurple}{RGB}{255,224,255}
\definecolor{lightgrey}{RGB}{77,77,77}

% Specify general colours
\colorlet{colorKeywordText}{CSviolet}
\colorlet{colorKeywordMargin}{CSgrey}
\colorlet{colorHighlight}{CSmagenta}
\colorlet{colorCite}{CScyan}
\colorlet{colorLink}{CSblue}

% Specify colours used for comments, questions, etc.
\colorlet{colorTodo}{darkblue}
\colorlet{colorTodoFill}{lightblue}
\colorlet{colorQuestion}{darkred}
\colorlet{colorQuestionFill}{lightred}
\colorlet{colorFeedback}{darkpurple}
\colorlet{colorFeedbackFill}{lightpurple}

% Specify colours used in environment boxes
\colorlet{colorBoxTheorem}{CSred}
\colorlet{colorBoxTheoremFill}{colorBoxTheorem!8!white}
\colorlet{colorBoxLemma}{CSred}
\colorlet{colorBoxCorollary}{CSred}
\colorlet{colorBoxProof}{CSlightgrey}
\colorlet{colorBoxProposition}{CSred}
\colorlet{colorBoxDefinition}{CSlightblue}
\colorlet{colorBoxExample}{CSgreen}
\colorlet{colorBoxExampleFill}{CSgreen!8!white}
\colorlet{colorBoxRemark}{CSgrey}
\colorlet{colorBoxRemarkFill}{CSgrey!8!white}
\colorlet{colorBoxProblem}{CSblack}
\colorlet{colorBoxConstraint}{CSviolet}
\colorlet{colorBoxConstraintFill}{CSviolet!8!white}

\colorlet{colorBoxRestate}{CSred}
\colorlet{colorBoxAssumption}{CSviolet}
\colorlet{colorBoxAssumptionFill}{CSviolet!8!white}
\colorlet{colorBoxSummary}{CSorange}
\colorlet{colorBoxSummaryFill}{CSorange!8!white}
\colorlet{colorBoxQuoteFill}{CSlightgrey!16!white}

 % Set used colours


%%%%%%%%%%%
% Debug
%%%%%%%%%%%
\iftoggle{debug}{
    \overfullrule=2cm % Show overfull hboxes with black line in margin
}{}


%%%%%%%%%%%
% Layout
%%%%%%%%%%%
\patchcmd{\part}{\thispagestyle{plain}}{\thispagestyle{empty}}{}{}

% Header
\clearpairofpagestyles
\ihead{\headmark}
\ohead*{\pagemark}

\makeatletter
\newcommand{\l@summary}[2]{%
    \begingroup%
        \let\@dotsep\@M%
        \csname l@#2\endcsname{#1}{}%
    \endgroup%
}%
\newcommand{\desctotoc}[1]{%
    \addtocontents{toc}{%
        \unexpanded{%
            \unexpanded{%
                {\small#1\par\medskip}%
            }%
        }%
    }%
}%

% Chapter and subtitle
\let\LaTeXStandardChapter\chapter%

\NewDocumentCommand{\chaptersubtitlefont}{+m}{%
    {\large \bfseries #1}%
}%

\RenewDocumentCommand{\chapter}{s+o+m+o}{%
    \IfBooleanTF{#1}{%
        \LaTeXStandardChapter*{#3}%
    }{%
        \IfNoValueTF{#4}{%
            \IfNoValueTF{#2}{%
                \LaTeXStandardChapter{#3}%
            }{%
                \LaTeXStandardChapter[#3]{#3}%
            }%
        }{%
            \IfNoValueTF{#2}{%
                \LaTeXStandardChapter[#3]{#3\\ \chaptersubtitlefont{#4}}%
            }{%
                \LaTeXStandardChapter[#3]{#3\\ \chaptersubtitlefont{#4}}%
            }%
        }%
    }%
    \IfValueTF{#4}{%
        \desctotoc{#4}%
    }{%
        %\desctotoc{}%
    }%
}%

\makeatother


% Fonts for headings
\addtokomafont{disposition}{\rmfamily\mdseries}
\addtokomafont{chapterprefix}{\large}
\setkomafont{dictumtext}{\normalfont\normalcolor\itshape\setlength{\parindent}{1em}\noindent}
\addtokomafont{subsubsection}{\normalfont\bfseries}% change font settings for paragraph heading
\addtokomafont{paragraph}{\normalfont\bfseries}
\RedeclareSectionCommand[afterskip=-.5em]{paragraph}% change horizontal skip after paragraph heading
\renewcommand{\sectioncatchphraseformat}[4]{%
    \Ifstr{#1}{paragraph}
    {\hskip #2#3#4\hskip .5em{\normalfont\textbar}}% new definition for paragraph heading
    {\hskip #2#3#4}% original definition for other levels like subparagraph
}

% Chapter thumb
\AddLayersToPageStyle{@everystyle@}{chapterthumb}
\addtokomafont{chapterthumb}{\bfseries}

% List indent
\setlist[itemize,1]{leftmargin=.175in}

% Tables
\setlength{\tabcolsep}{3pt}
\setlength{\arraycolsep}{2pt}

% Hyperref
\hypersetup{
    colorlinks=true,
    linkcolor=colorLink,
    citecolor=colorCite,
    pdftitle=\TITLE{},
    pdfauthor=\AUTHOR{}
}

% Quotes
% Add quotation marks for environments "displayquote" and related:
\renewcommand{\mkbegdispquote}[2]{\openautoquote}
\renewcommand{\mkenddispquote}[2]{\closeautoquote#1#2}



%%%%%%%%%%%
% Bibliography
%%%%%%%%%%%

\renewcommand*{\citesetup}{%
    \biburlsetup
    \frenchspacing
}

% Define two different bibliographies (general and own publications) with different labels
\DeclareFieldFormat{labelnumberwidth}{\mkbibbrackets{#1}}
\renewbibmacro*{cite}{%
    \printtext[bibhyperref]{%
        \printfield{labelprefix}%
        \ifkeyword{own}%
        {%
            \printfield{labelnumber}%
        }{%
            \printfield{labelalpha}%
            \printfield{extraalpha}%
        }%
    }%
}

% Bibliography with numbered labels
\defbibenvironment{bibliographyNUM}
{%
    \list%
    {%
        \printtext[labelnumberwidth]{%
            \printfield{prefixnumber}%
            \printfield{labelnumber}%
        }%
    }{%
    \setlength{\labelwidth}{\labelnumberwidth}%
    \setlength{\leftmargin}{\labelwidth}%
    %\addtolength{\leftmargin}{\labelsep}%
    \setlength{\labelsep}{\biblabelsep}%
    \setlength{\itemsep}{\bibitemsep}%
    \setlength{\parsep}{\bibparsep}%
    }%
    \renewcommand*{\makelabel}[1]{\hss##1}%
}{\endlist}{\item}

% Bibliography with non-numbered labels
\defbibenvironment{bibliographyNONUM}%
{%
    \list%
    {}%
    {%
        \setlength{\leftmargin}{\bibhang}%
        \setlength{\itemindent}{-\leftmargin}%
        \setlength{\itemsep}{\bibitemsep}%
        \setlength{\parsep}{\bibparsep}%
    }%
}{\endlist}{\item}

% Source of bibliography
% Papers from oneself are given the additional tag 'own', other paper are tagged 'other'
% TODO: Change to own surname
\DeclareSourcemap{
    \maps[datatype=bibtex]{
    \map[overwrite]{
        \step[fieldsource=author, notmatch=\regexp{Knuth}, final]
        \step[fieldsource=keywords, match=\regexp{\A.+\Z}, final]
        \step[fieldset=keywords, fieldvalue={,other}, append]
    }
    \map{
        \step[fieldsource=author, notmatch=\regexp{(Knuth)}, final]
        \step[fieldset=keywords, fieldvalue=other, append]
    }
    \map[overwrite]{
        \step[fieldsource=author, match=\regexp{(Knuth)}, final]
        \step[fieldsource=keywords, match=\regexp{\A.+\Z}, final]
        \step[fieldset=keywords, fieldvalue={,own}, append]
    }
    \map{
        \step[fieldsource=author, match=\regexp{(Knuth)}, final]
        \step[fieldset=keywords, fieldvalue=own]
    }
  }
}

\defbibheading{subbibliography}[\refname]{\section*{#1}}
\DeclareBibliographyCategory{cited}
\AtEveryCitekey{\addtocategory{cited}{\thefield{entrykey}}}

\assignrefcontextkeyws[sorting=ynt]{own}

\makeatletter
% Full citation
\DeclareCiteCommand{\bfullcite}
{\usebibmacro{prenote}}
{
    \renewcommand*{\newunitpunct}{\addcomma\space}%
    \defcounter{maxnames}{99}%
    \defcounter{minnames}{99}%
    \settoggle{blx@useauthor}{false}%
    \settoggle{blx@useeditor}{false}%
    \settoggle{blx@usetranslator}{false}%
    \usedriver%
    {%
        \DeclareNameAlias{sortname}{default}%
    }{%
        \thefield{entrytype}%
    }%
}%
{\multicitedelim}%
{\usebibmacro{postnote}}%
\makeatother

% Bibliography style
\renewcommand*{\labelalphaothers}{\textsuperscript{+}}%
\renewbibmacro{in:}{}%
\DeclareFieldFormat*{title}{\mkbibquote{\textsl{#1}\isdot}}%
\DeclareFieldFormat*{year}{#1\isdot}%
\AtEveryBibitem{% Clean up the bibtex rather than editing it
    \clearlist{address}%
    \clearlist{location}%
    \clearfield{isbn}%
    \clearfield{issn}%
}%



%%%%%%%%%%%
% Theorem settings
%%%%%%%%%%%

%% theorem settings
\theorembodyfont{}      
%\theoremsymbol{\ensuremath{\triangle}}
\theoremprework{}
\theorempostwork{}
\theoremseparator{\newline}

\newtheorem{thmcounter}{thmcounter}[chapter]
\newtheorem{deficounter}{deficounter}[chapter]
\newtheorem{excounter}{excounter}[chapter]
\newtheorem{remcounter}{remcounter}[chapter]

% Theorems, lemmas, ...
\theoremstyle{plain}
\newtheorem{theorem}[thmcounter]{Theorem}
\newtheorem{lemma}[thmcounter]{Lemma}
\newtheorem{corollary}[thmcounter]{Corollary}
\newtheorem{proposition}[thmcounter]{Proposition}

% Definitions, remarks, examples, ...
\theorembodyfont{\normalfont}
\newtheorem{definition}[deficounter]{Definition}
\newtheorem{example}[excounter]{Example}
\newtheorem{remark}[remcounter]{Remark}
\newtheorem{problem}{Problem}[chapter]

% Custom constraints
\newcounter{constraintcounter}
\newtheorem{constraint}[constraintcounter]{Constraint}
\crefname{constraint}{Constraint}{Constraints} % Define naming for cleveref
% Define custom environment which automatically sets labelname for correct \nameref
\makeatletter
\newenvironment{custom_constraint}[1]{\begin{constraint}[#1]\let\@currentlabelname\constraintname}{\end{constraint}}
\makeatother

% Proofs
\theoremstyle{nonumberplain}
\theoremheaderfont{\normalfont\itshape}
\theorembodyfont{\normalfont}
\theoremseparator{.\,}
\theoremsymbol{\ensuremath{\blacksquare}}
\newtheorem{proof}{Proof}
\newtheorem{proofsketch}{Proof (sketch)}


%%%%%%%%%%%
% Environment boxes
%%%%%%%%%%%

% Set general styles
\tcbset{
    colorboxbase/.style={
        breakable, % allow multi page boxes
        boxsep=0mm,
        left=3mm,
        before skip=0.5\baselineskip,
        after skip=0.5\baselineskip,
        after={}
    },
    boxbackground/.style={
        enhanced,
        colorboxbase,
        top=2mm,
        bottom=2mm,
        right=2mm,
        frame hidden % disable borders
    },
    boxline/.style={
        blanker,
        colorboxbase,
        top=0.5mm,
        bottom=0.5mm,
        right=0mm,
        frame hidden, % disable borders
        colback=white
    }
}

% Define styles specific to environment boxes
\tcolorboxenvironment{theorem}{
    boxbackground,
    colback=colorBoxTheoremFill,
    borderline west={0.6mm}{0pt}{colorBoxTheorem},
}
\tcolorboxenvironment{lemma}{
    boxline,
    borderline west={0.6mm}{0pt}{colorBoxLemma},
}
\tcolorboxenvironment{corollary}{
    boxline,
    borderline west={0.6mm}{0pt}{colorBoxCorollary},
}
\tcolorboxenvironment{proposition}{
    boxline,
    borderline west={0.6mm}{0pt}{colorBoxProposition},
}
\tcolorboxenvironment{proof}{
    boxline,
    borderline west={0.6mm}{0pt}{colorBoxProof},
}
\tcolorboxenvironment{proofsketch}{
    boxline,
    borderline west={0.6mm}{0pt}{colorBoxProof},
}

\tcolorboxenvironment{definition}{
    boxline,
    borderline west={0.6mm}{0pt}{colorBoxDefinition},
}
\tcolorboxenvironment{example}{
    boxbackground,
    colback=colorBoxExampleFill,
    borderline west={0.6mm}{0pt}{colorBoxExample},
}
\tcolorboxenvironment{remark}{
    boxbackground,
    colback=colorBoxRemarkFill,
    borderline west={0.6mm}{0pt}{colorBoxRemark},
}
\tcolorboxenvironment{problem}{
    boxline,
    borderline west={0.6mm}{0pt}{colorBoxProblem},
}
\tcolorboxenvironment{custom_constraint}{
    boxbackground,
    colback=colorBoxConstraintFill,
    borderline west={0.6mm}{0pt}{colorBoxConstraint},
}

\tcolorboxenvironment{displayquote}{
    boxbackground,
    colback=colorBoxQuoteFill,
}

%%%%%%%%%%%
% More boxes
%%%%%%%%%%%

% Box for restating theorems/lemmas/... in the appendix
\newtcolorbox{restatebox}{
    boxline,
    borderline west={0.6mm}{0pt}{colorBoxRestate},
}

% Box for assumptions in the beginning of the chapter
\newtcolorbox{assumptionbox}{
    colback=colorBoxAssumptionFill,
    colframe=colorBoxAssumption,
    fonttitle=\bfseries,
    title=Assumptions
}

% Box for important lessons at the end of the chapter
\newtcolorbox{summarybox}{
    colback=colorBoxSummaryFill,
    colframe=colorBoxSummary,
    fonttitle=\bfseries,
    title=Summary
}

% Define summary as itemization in summary box
\newenvironment{summary}{%
    \phantomsection% Needed for correct links with hyperref
    \addcontentsline{toc}{section}{Summary}% Add into TOC
    \begin{summarybox}\begin{itemize}[label=\ding{220}]}
{\end{itemize}\end{summarybox}}


%%%%%%%%%%%
% Figures
%%%%%%%%%%%

% \vspace between stacked subfigures
\newcommand*{\subvspace}{\vspace{0.5cm}}
